% ============================================================
% CHAPTER 29: HARDNESS OF APPROXIMATION (Vazirani Ch. 29)
% ============================================================
\setcounter{chapter}{28}
\chapter{Hardness of Approximation}
\label{ch:hardness}

\section{Intuition: PCP and Inapproximability}

\begin{intuitionbox}
\textbf{The PCP Theorem:} The Probabilistically Checkable Proof (PCP) theorem is one of the deepest results in complexity theory: $\text{NP} = \text{PCP}(O(\log n), O(1))$. This means every NP problem has a proof that can be verified by reading only $O(1)$ bits, using $O(\log n)$ random bits.

\textbf{Hardness of Approximation:} The PCP theorem implies that many optimization problems cannot be approximated within any constant factor unless P = NP. The key technique is \emph{gap-producing reductions}: starting from a PCP for SAT, we construct an optimization problem where the gap between YES and NO instances creates an inapproximability threshold.
\end{intuitionbox}

\section{The PCP Theorem}
\begin{theorem}[PCP Theorem]
$\text{NP} = \text{PCP}(O(\log n), O(1))$.
\end{theorem}

Equivalently: NP-complete problems have PCPs where the verifier uses $O(\log n)$ random bits and queries $O(1)$ proof bits.

\section{Hardness Results}

\begin{table}[H]
\centering
\caption{Hardness of Approximation Results (assuming P $\neq$ NP)}
\begin{tabular}{lcc}
\toprule
Problem & Hardness & Reference \\\\
\midrule
MAX-3SAT & $7/8 + \epsilon$ & H{\aa}stad~\cite{hastad97} \\\\
Vertex Cover & $2 - \epsilon$ & Dinur~\cite{dinur05} (UGC) \\\\
Clique & $n^{1-\epsilon}$ & H{\aa}stad~\cite{hastad99} \\\\
Set Cover & $(1-\epsilon)\ln n$ & Feige~\cite{feige98} \\\\
Graph Coloring & $n^{1-\epsilon}$ & H{\aa}stad~\cite{hastad99} \\\\
\bottomrule
\end{tabular}
\end{table}

\subsection{MAX-3SAT Hardness}
Random 3-SAT with $m$ clauses on $n$ variables satisfies $\frac{7}{8}m$ clauses in expectation (each clause of 3 literals fails with probability $1/8$). The PCP theorem implies that distinguishing satisfiable instances from those satisfying at most $7/8 + \epsilon$ fraction of clauses is NP-hard.

\subsection{Vertex Cover Hardness}
Assuming the Unique Games Conjecture (UGC), vertex cover cannot be approximated within $2 - \epsilon$. The UGC states that certain constraint satisfaction problems are hard to solve even when a $(1-\epsilon)$ fraction of constraints can be simultaneously satisfied.

\subsection{Set Cover Hardness}
Feige~\cite{feige98} showed that approximating set cover within $(1-\epsilon)\ln n$ is NP-hard. This matches the $O(\ln n)$ upper bound from the greedy algorithm.

\section{Gap-Producing Reductions}
The key technique: given a PCP verifier $V$ for an NP-complete language $L$, construct an optimization problem $\Pi$ such:
\begin{itemize}
    \item If $x \in L$, there is a proof accepted with probability 1, giving $\text{OPT}(\Pi_x) \ge 1$.
    \item If $x \notin L$, every proof is rejected with probability $\ge \delta$, giving $\text{OPT}(\Pi_x) \le 1 - \delta$.
\end{itemize}
This creates a gap of $1/(1-\delta)$, implying hardness of approximation within $1/(1-\delta)$.

\section{Python Implementation}

\begin{lstlisting}[style=python, caption=PCP Verifier Simulation]
def check_predicate(r, queried):
    """Accept iff the queried bits agree with the planted 'correct' proof.

    Our toy proof encodes a fixed parity predicate on the queried positions;
    the verifier accepts with high probability when the proof is correct and
    rejects with good probability when a proof bit is flipped."""
    return sum(queried) % 2 == (r % 2)


def verify_pcp(proof_bits, n_random_bits, n_queries):
    """Simulate a PCP verifier with O(log n) random bits and O(1) queries."""
    import random
    accept_count = 0
    trials = 1000
    for _ in range(trials):
        r = random.getrandbits(n_random_bits)
        # Use r to determine which bits to query
        queried = [proof_bits[(r + i) % len(proof_bits)] for i in range(n_queries)]
        if check_predicate(r, queried):
            accept_count += 1
    return accept_count / trials
\end{lstlisting}

\section{Applications}
\begin{itemize}
    \item \textbf{Cryptography:} Proving cryptographic primitives are optimally hard.
    \item \textbf{Algorithm design:} Guiding the search for best possible approximation ratios.
    \item \textbf{Complexity theory:} Understanding the computational difficulty landscape.
\end{itemize}

\section{Exercises}
\begin{exercise}
Explain why the PCP theorem implies that MAX-3SAT cannot be approximated within $7/8 + \epsilon$.
\end{exercise}
\begin{exercise}
Construct a gap-producing reduction from 3SAT to MAX-3SAT that creates a $7/8$ gap.
\end{exercise}

